%=====================================================================
% Implementación · Restricciones
%   Las tablas de CHECK y de valores por omisión las genera
%   bd/generar_tablas.ps1 a partir de bd/crear_tablas.sql.
%=====================================================================
\subsection*{Restricciones}
\addcontentsline{toc}{subsection}{Restricciones}

Además de las llaves, el script declara las reglas que la base puede hacer cumplir por sí misma, para que no dependan de que el servidor de partidas las recuerde. Las cifras de esta subsección se cuentan sobre \at{crear\_tablas.sql} al generar el documento.

\subsubsection*{\texttt{NOT NULL}}
\addcontentsline{toc}{subsubsection}{NOT NULL}

\bdNumNoNulas{} de las \bdNumColumnas{} columnas son \at{NOT NULL}; \bdNumCalculadas{} más son calculadas y no admiten la restricción. Una columna admite nulo solo por una de tres razones, y la columna \emph{Nulo} de las tablas de atributos lo indica en cada caso:

\begin{reglas}
  \item \textbf{Aplica a un subtipo.} El correo y la contraseña, que no tiene el invitado; la fecha de fin, que no tiene una sanción permanente; las casillas de una jugada, que no tiene un muro. Un \at{CHECK} con el discriminador exige el valor en el subtipo al que aplica y lo prohíbe en el otro.
  \item \textbf{Todavía no ocurre.} El resultado de una participación en curso, la fecha de fin de una sesión abierta, el moderador de un reporte pendiente. Un \at{CHECK} ata la columna al estado de la fila.
  \item \textbf{Es opcional.} La descripción de un reporte, la partida que adjunta, el nombre del dispositivo.
\end{reglas}

\subsubsection*{\texttt{UNIQUE}}
\addcontentsline{toc}{subsubsection}{UNIQUE}

\bdNumUnicas{} restricciones \at{UNIQUE} declaran las otras llaves candidatas, como el \at{nickname} o el \at{token\_hash} de una sesión. Hay reglas de unicidad que solo valen para parte de las filas: un correo único entre las cuentas que lo tienen, un solo token vigente por cuenta, un código de sala único entre las salas no cerradas. Una restricción \at{UNIQUE} no puede expresarlas, porque trata todas las filas por igual y en SQL Server admite un solo nulo; por eso se declaran como \bdNumIndicesUnicos{} índices únicos filtrados, \at{CREATE UNIQUE INDEX ... WHERE ...}. Las dos clases están en la tabla \textit{Restricciones de unicidad} del modelo de datos.

\subsubsection*{\texttt{CHECK}}
\addcontentsline{toc}{subsubsection}{CHECK}

Las \bdNumChecks{} restricciones \at{CHECK} son de cuatro clases:

\begin{reglas}
  \item \textbf{Dominios cerrados.} Cada estado, tipo, ámbito u origen admite solo los valores de la tabla de dominios del análisis CRUD, por ejemplo \at{estado\_cuenta IN ('PENDIENTE', 'ACTIVA', ...)}. Los catálogos que pueden crecer, como las ranuras, no llevan \at{CHECK}: añadir un valor es insertar una fila.
  \item \textbf{Rangos.} Contadores no negativos, el icono entre 1 y 32, de 1 a 20 muros por jugador, relojes de 3, 5 o 10 minutos, probabilidades entre 0 y 100.
  \item \textbf{Coherencia entre columnas.} Qué columnas aplican a cada valor de un discriminador o de un estado: una partida contra la IA tiene nivel y opciones y las demás no; un reporte en revisión tiene moderador y un reporte pendiente no.
  \item \textbf{Reglas de la fila.} Que nadie se reporte, se invite o se bloquee a sí mismo, que una amistad se guarde una sola vez por par ordenado, o que el elo máximo no sea menor que el actual.
\end{reglas}

% Archivo generado por bd/generar_tablas.ps1 a partir de bd/crear_tablas.sql. No editar a mano.
\begin{center}\small
\begin{longtable}{|L{0.19\textwidth}|L{0.23\textwidth}|L{0.48\textwidth}|}
\hline
\textbf{Entidad} & \textbf{Restricción} & \textbf{Condición} \\ \hline
\endfirsthead
\hline
\textbf{Entidad} & \textbf{Restricción} & \textbf{Condición} \\ \hline
\endhead
\multicolumn{3}{|l|}{\textbf{Catálogos precargados (\mbox{D-09})}} \\ \hline
\textit{Objeto\-Cosmetico} & \texttt{CK\_\allowbreak{}Objeto\allowbreak{}Cosmetico\_\allowbreak{}rareza} & \texttt{rareza IN (\textquotesingle{}COMUN\textquotesingle{},\allowbreak{}\textquotesingle{}RARO\textquotesingle{},\allowbreak{}\textquotesingle{}EPICO\textquotesingle{},\allowbreak{}\textquotesingle{}LEGENDARIO\textquotesingle{})} \\ \hline
\textit{Objeto\-Cosmetico} & \texttt{CK\_\allowbreak{}Objeto\allowbreak{}Cosmetico\_\allowbreak{}rangos} & \texttt{precio >= 0 AND nivel\_\allowbreak{}requerido >= 1} \\ \hline
\textit{Objeto\-Cosmetico} & \texttt{CK\_\allowbreak{}Objeto\allowbreak{}Cosmetico\_\allowbreak{}predeterminado} & \texttt{es\_\allowbreak{}predeterminado = 0 OR es\_\allowbreak{}inicial = 1} \\ \hline
\textit{Objeto\-Cosmetico} & \texttt{CK\_\allowbreak{}Objeto\allowbreak{}Cosmetico\_\allowbreak{}venta} & \texttt{a\_\allowbreak{}la\_\allowbreak{}venta = 0 OR activo = 1} \\ \hline
\textit{Modo} & \texttt{CK\_\allowbreak{}Modo\_\allowbreak{}parametros} & \texttt{tamano\_\allowbreak{}tablero IN (7,\allowbreak{} 9) AND num\_\allowbreak{}jugadores IN (2,\allowbreak{} 4) AND muros\_\allowbreak{}por\_\allowbreak{}jugador BETWEEN 1 AND 20} \\ \hline
\textit{Division} & \texttt{CK\_\allowbreak{}Division\_\allowbreak{}umbral} & \texttt{elo\_\allowbreak{}descenso <= elo\_\allowbreak{}minimo} \\ \hline
\textit{Leccion} & \texttt{CK\_\allowbreak{}Leccion\_\allowbreak{}pasos} & \texttt{num\_\allowbreak{}pasos > 0} \\ \hline
\textit{Tipo\-Caja} & \texttt{CK\_\allowbreak{}Tipo\allowbreak{}Caja\_\allowbreak{}precio} & \texttt{precio >= 0} \\ \hline
\textit{Tipo\-Caja\-Objeto} & \texttt{CK\_\allowbreak{}Tipo\allowbreak{}Caja\allowbreak{}Objeto\_\allowbreak{}probabilidad} & \texttt{probabilidad > 0 AND probabilidad <= 100} \\ \hline
\textit{Palabra\-Prohibida} & \texttt{CK\_\allowbreak{}Palabra\allowbreak{}Prohibida\_\allowbreak{}idioma} & \texttt{idioma IN (\textquotesingle{}es-MX\textquotesingle{},\allowbreak{}\textquotesingle{}en\textquotesingle{})} \\ \hline
\textit{Palabra\-Prohibida} & \texttt{CK\_\allowbreak{}Palabra\allowbreak{}Prohibida\_\allowbreak{}ambito} & \texttt{ambito IN (\textquotesingle{}NICKNAME\textquotesingle{},\allowbreak{}\textquotesingle{}CHAT\textquotesingle{},\allowbreak{}\textquotesingle{}AMBOS\textquotesingle{})} \\ \hline
\multicolumn{3}{|l|}{\textbf{Cuenta y acceso}} \\ \hline
\textit{Usuario} & \texttt{CK\_\allowbreak{}Usuario\_\allowbreak{}tipo\_\allowbreak{}cuenta} & \texttt{tipo\_\allowbreak{}cuenta IN (\textquotesingle{}INVITADO\textquotesingle{},\allowbreak{}\textquotesingle{}REGISTRADA\textquotesingle{})} \\ \hline
\textit{Usuario} & \texttt{CK\_\allowbreak{}Usuario\_\allowbreak{}estado\_\allowbreak{}cuenta} & \texttt{estado\_\allowbreak{}cuenta IN (\textquotesingle{}PENDIENTE\textquotesingle{},\allowbreak{}\textquotesingle{}ACTIVA\textquotesingle{},\allowbreak{}\textquotesingle{}SUSPENDIDA\textquotesingle{},\allowbreak{}\textquotesingle{}BANEADA\textquotesingle{},\allowbreak{}\textquotesingle{}ELIMINADA\textquotesingle{})} \\ \hline
\textit{Usuario} & \texttt{CK\_\allowbreak{}Usuario\_\allowbreak{}rol} & \texttt{rol IN (\textquotesingle{}JUGADOR\textquotesingle{},\allowbreak{}\textquotesingle{}MODERADOR\textquotesingle{},\allowbreak{}\textquotesingle{}ADMINISTRADOR\textquotesingle{})} \\ \hline
\textit{Usuario} & \texttt{CK\_\allowbreak{}Usuario\_\allowbreak{}idioma} & \texttt{idioma\_\allowbreak{}preferido IN (\textquotesingle{}es-MX\textquotesingle{},\allowbreak{}\textquotesingle{}en\textquotesingle{})} \\ \hline
\textit{Usuario} & \texttt{CK\_\allowbreak{}Usuario\_\allowbreak{}credenciales} & \texttt{(tipo\_\allowbreak{}cuenta = \textquotesingle{}INVITADO\textquotesingle{} AND correo IS NULL AND contrasena\_\allowbreak{}hash IS NULL AND contrasena\_\allowbreak{}sal IS NULL) OR (tipo\_\allowbreak{}cuenta = \textquotesingle{}REGISTRADA\textquotesingle{} AND correo IS NOT NULL)} \\ \hline
\textit{Usuario} & \texttt{CK\_\allowbreak{}Usuario\_\allowbreak{}rol\_\allowbreak{}registrada} & \texttt{rol = \textquotesingle{}JUGADOR\textquotesingle{} OR tipo\_\allowbreak{}cuenta = \textquotesingle{}REGISTRADA\textquotesingle{}} \\ \hline
\textit{Usuario} & \texttt{CK\_\allowbreak{}Usuario\_\allowbreak{}baja} & \texttt{(estado\_\allowbreak{}cuenta = \textquotesingle{}ELIMINADA\textquotesingle{} AND fecha\_\allowbreak{}baja IS NOT NULL) OR (estado\_\allowbreak{}cuenta <> \textquotesingle{}ELIMINADA\textquotesingle{} AND fecha\_\allowbreak{}baja IS NULL)} \\ \hline
\textit{Usuario} & \texttt{CK\_\allowbreak{}Usuario\_\allowbreak{}rangos} & \texttt{id\_\allowbreak{}icono BETWEEN 1 AND 32 AND nivel >= 1 AND experiencia >= 0 AND saldo\_\allowbreak{}monedas >= 0 AND cajas\_\allowbreak{}sin\_\allowbreak{}raro BETWEEN 0 AND 10} \\ \hline
\textit{Sesion} & \texttt{CK\_\allowbreak{}Sesion\_\allowbreak{}motivo} & \texttt{motivo\_\allowbreak{}cierre IN (\textquotesingle{}CIERRE\_\allowbreak{}VOLUNTARIO\textquotesingle{},\allowbreak{}\textquotesingle{}CIERRE\_\allowbreak{}REMOTO\textquotesingle{},\allowbreak{}\textquotesingle{}CAMBIO\_\allowbreak{}CREDENCIALES\textquotesingle{},\allowbreak{}\textquotesingle{}BAJA\_\allowbreak{}CUENTA\textquotesingle{},\allowbreak{}\textquotesingle{}SANCION\textquotesingle{})} \\ \hline
\textit{Sesion} & \texttt{CK\_\allowbreak{}Sesion\_\allowbreak{}cierre} & \texttt{(fecha\_\allowbreak{}fin IS NULL AND motivo\_\allowbreak{}cierre IS NULL) OR (fecha\_\allowbreak{}fin IS NOT NULL AND motivo\_\allowbreak{}cierre IS NOT NULL)} \\ \hline
\textit{Sesion} & \texttt{CK\_\allowbreak{}Sesion\_\allowbreak{}vigencia} & \texttt{fecha\_\allowbreak{}expiracion > fecha\_\allowbreak{}inicio} \\ \hline
\textit{Codigo\-Segundo\-Factor} & \texttt{CK\_\allowbreak{}Codigo\allowbreak{}Segundo\allowbreak{}Factor\_\allowbreak{}intentos} & \texttt{intentos BETWEEN 0 AND 3} \\ \hline
\textit{Codigo\-Segundo\-Factor} & \texttt{CK\_\allowbreak{}Codigo\allowbreak{}Segundo\allowbreak{}Factor\_\allowbreak{}estado} & \texttt{estado IN (\textquotesingle{}PENDIENTE\textquotesingle{},\allowbreak{}\textquotesingle{}CONSUMIDO\textquotesingle{},\allowbreak{}\textquotesingle{}INVALIDADO\textquotesingle{})} \\ \hline
\textit{Token\-Verificacion\-Correo} & \texttt{CK\_\allowbreak{}Token\allowbreak{}Verificacion\allowbreak{}Correo\_\allowbreak{}proposito} & \texttt{(proposito = \textquotesingle{}ALTA\textquotesingle{} AND correo\_\allowbreak{}destino IS NULL) OR (proposito = \textquotesingle{}CAMBIO\_\allowbreak{}CORREO\textquotesingle{} AND correo\_\allowbreak{}destino IS NOT NULL)} \\ \hline
\textit{Token\-Verificacion\-Correo} & \texttt{CK\_\allowbreak{}Token\allowbreak{}Verificacion\allowbreak{}Correo\_\allowbreak{}estado} & \texttt{estado IN (\textquotesingle{}PENDIENTE\textquotesingle{},\allowbreak{}\textquotesingle{}USADO\textquotesingle{},\allowbreak{}\textquotesingle{}INVALIDADO\textquotesingle{})} \\ \hline
\textit{Token\-Recuperacion} & \texttt{CK\_\allowbreak{}Token\allowbreak{}Recuperacion\_\allowbreak{}intentos} & \texttt{intentos BETWEEN 0 AND 3} \\ \hline
\textit{Token\-Recuperacion} & \texttt{CK\_\allowbreak{}Token\allowbreak{}Recuperacion\_\allowbreak{}estado} & \texttt{estado IN (\textquotesingle{}PENDIENTE\textquotesingle{},\allowbreak{}\textquotesingle{}USADO\textquotesingle{},\allowbreak{}\textquotesingle{}INVALIDADO\textquotesingle{})} \\ \hline
\textit{Aceptacion\-Terminos} & \texttt{CK\_\allowbreak{}Aceptacion\allowbreak{}Terminos\_\allowbreak{}idioma} & \texttt{idioma IN (\textquotesingle{}es-MX\textquotesingle{},\allowbreak{}\textquotesingle{}en\textquotesingle{})} \\ \hline
\multicolumn{3}{|l|}{\textbf{Perfil y personalización}} \\ \hline
\textit{Avatar} & \texttt{CK\_\allowbreak{}Avatar\_\allowbreak{}formato} & \texttt{formato IN (\textquotesingle{}JPG\textquotesingle{},\allowbreak{}\textquotesingle{}PNG\textquotesingle{})} \\ \hline
\textit{Avatar} & \texttt{CK\_\allowbreak{}Avatar\_\allowbreak{}tamano} & \texttt{tamano\_\allowbreak{}bytes BETWEEN 1 AND 2097152} \\ \hline
\textit{Enlace\-Red} & \texttt{CK\_\allowbreak{}Enlace\allowbreak{}Red\_\allowbreak{}orden} & \texttt{orden BETWEEN 1 AND 4} \\ \hline
\textit{Usuario\-Objeto} & \texttt{CK\_\allowbreak{}Usuario\allowbreak{}Objeto\_\allowbreak{}origen} & \texttt{origen IN (\textquotesingle{}INICIAL\textquotesingle{},\allowbreak{}\textquotesingle{}COMPRA\textquotesingle{},\allowbreak{}\textquotesingle{}CAJA\textquotesingle{},\allowbreak{}\textquotesingle{}NIVEL\textquotesingle{})} \\ \hline
\multicolumn{3}{|l|}{\textbf{Partida}} \\ \hline
\textit{Partida} & \texttt{CK\_\allowbreak{}Partida\_\allowbreak{}tipo} & \texttt{tipo IN (\textquotesingle{}CLASIFICATORIA\textquotesingle{},\allowbreak{}\textquotesingle{}PRIVADA\textquotesingle{},\allowbreak{}\textquotesingle{}IA\textquotesingle{})} \\ \hline
\textit{Partida} & \texttt{CK\_\allowbreak{}Partida\_\allowbreak{}estado} & \texttt{estado IN (\textquotesingle{}EN\_\allowbreak{}CURSO\textquotesingle{},\allowbreak{}\textquotesingle{}PENDIENTE\_\allowbreak{}DE\_\allowbreak{}CIERRE\textquotesingle{},\allowbreak{}\textquotesingle{}FINALIZADA\textquotesingle{})} \\ \hline
\textit{Partida} & \texttt{CK\_\allowbreak{}Partida\_\allowbreak{}forma} & \texttt{forma\_\allowbreak{}termino IN (\textquotesingle{}META\textquotesingle{},\allowbreak{}\textquotesingle{}TIEMPO\_\allowbreak{}AGOTADO\textquotesingle{},\allowbreak{}\textquotesingle{}RENDICION\textquotesingle{},\allowbreak{}\textquotesingle{}TABLAS\textquotesingle{},\allowbreak{}\textquotesingle{}ABANDONO\textquotesingle{})} \\ \hline
\textit{Partida} & \texttt{CK\_\allowbreak{}Partida\_\allowbreak{}reloj} & \texttt{minutos\_\allowbreak{}reloj IN (3,\allowbreak{} 5,\allowbreak{} 10)} \\ \hline
\textit{Partida} & \texttt{CK\_\allowbreak{}Partida\_\allowbreak{}ia} & \texttt{(tipo = \textquotesingle{}IA\textquotesingle{} AND id\_\allowbreak{}nivel\_\allowbreak{}ia IS NOT NULL AND permite\_\allowbreak{}deshacer IS NOT NULL AND permite\_\allowbreak{}sugerencias IS NOT NULL AND deshacer\_\allowbreak{}usados IS NOT NULL AND sugerencias\_\allowbreak{}usadas IS NOT NULL) OR (tipo <> \textquotesingle{}IA\textquotesingle{} AND id\_\allowbreak{}nivel\_\allowbreak{}ia IS NULL AND permite\_\allowbreak{}deshacer IS NULL AND permite\_\allowbreak{}sugerencias IS NULL AND deshacer\_\allowbreak{}usados IS NULL AND sugerencias\_\allowbreak{}usadas IS NULL AND minutos\_\allowbreak{}reloj IS NOT NULL)} \\ \hline
\textit{Partida} & \texttt{CK\_\allowbreak{}Partida\_\allowbreak{}cierre} & \texttt{(estado = \textquotesingle{}EN\_\allowbreak{}CURSO\textquotesingle{} AND fecha\_\allowbreak{}fin IS NULL AND forma\_\allowbreak{}termino IS NULL) OR (estado = \textquotesingle{}PENDIENTE\_\allowbreak{}DE\_\allowbreak{}CIERRE\textquotesingle{} AND fecha\_\allowbreak{}fin IS NULL AND forma\_\allowbreak{}termino IS NOT NULL) OR (estado = \textquotesingle{}FINALIZADA\textquotesingle{} AND fecha\_\allowbreak{}fin IS NOT NULL AND forma\_\allowbreak{}termino IS NOT NULL)} \\ \hline
\textit{Participacion} & \texttt{CK\_\allowbreak{}Participacion\_\allowbreak{}resultado} & \texttt{resultado IN (\textquotesingle{}GANADA\textquotesingle{},\allowbreak{}\textquotesingle{}PERDIDA\textquotesingle{},\allowbreak{}\textquotesingle{}TABLAS\textquotesingle{})} \\ \hline
\textit{Participacion} & \texttt{CK\_\allowbreak{}Participacion\_\allowbreak{}forma} & \texttt{forma\_\allowbreak{}termino IN (\textquotesingle{}META\textquotesingle{},\allowbreak{}\textquotesingle{}TIEMPO\_\allowbreak{}AGOTADO\textquotesingle{},\allowbreak{}\textquotesingle{}RENDICION\textquotesingle{},\allowbreak{}\textquotesingle{}ABANDONO\textquotesingle{})} \\ \hline
\textit{Participacion} & \texttt{CK\_\allowbreak{}Participacion\_\allowbreak{}rangos} & \texttt{orden\_\allowbreak{}turno BETWEEN 1 AND 4 AND posicion BETWEEN 1 AND 4 AND muros\_\allowbreak{}restantes <= 20} \\ \hline
\textit{Participacion} & \texttt{CK\_\allowbreak{}Participacion\_\allowbreak{}conexion} & \texttt{conectado = 1 OR fecha\_\allowbreak{}desconexion IS NOT NULL} \\ \hline
\textit{Jugada} & \texttt{CK\_\allowbreak{}Jugada\_\allowbreak{}tipo} & \texttt{(tipo = \textquotesingle{}MOVIMIENTO\textquotesingle{} AND casilla\_\allowbreak{}origen IS NOT NULL AND casilla\_\allowbreak{}destino IS NOT NULL AND surco IS NULL AND orientacion IS NULL) OR (tipo = \textquotesingle{}MURO\textquotesingle{} AND surco IS NOT NULL AND orientacion IN (\textquotesingle{}HORIZONTAL\textquotesingle{},\allowbreak{}\textquotesingle{}VERTICAL\textquotesingle{}) AND casilla\_\allowbreak{}origen IS NULL AND casilla\_\allowbreak{}destino IS NULL)} \\ \hline
\textit{Jugada} & \texttt{CK\_\allowbreak{}Jugada\_\allowbreak{}rangos} & \texttt{numero\_\allowbreak{}jugada >= 1 AND tiempo\_\allowbreak{}consumido >= 0} \\ \hline
\textit{Oferta\-Tablas} & \texttt{CK\_\allowbreak{}Oferta\allowbreak{}Tablas\_\allowbreak{}estado} & \texttt{estado IN (\textquotesingle{}PENDIENTE\textquotesingle{},\allowbreak{}\textquotesingle{}ACEPTADA\textquotesingle{},\allowbreak{}\textquotesingle{}RECHAZADA\textquotesingle{},\allowbreak{}\textquotesingle{}CADUCADA\textquotesingle{})} \\ \hline
\textit{Oferta\-Tablas} & \texttt{CK\_\allowbreak{}Oferta\allowbreak{}Tablas\_\allowbreak{}respuesta} & \texttt{(estado = \textquotesingle{}PENDIENTE\textquotesingle{} AND fecha\_\allowbreak{}respuesta IS NULL) OR (estado <> \textquotesingle{}PENDIENTE\textquotesingle{} AND fecha\_\allowbreak{}respuesta IS NOT NULL)} \\ \hline
\multicolumn{3}{|l|}{\textbf{Emparejamiento, salas y chat}} \\ \hline
\textit{Cola\-Emparejamiento} & \texttt{CK\_\allowbreak{}Cola\allowbreak{}Emparejamiento\_\allowbreak{}reloj} & \texttt{minutos\_\allowbreak{}reloj IN (3,\allowbreak{} 5,\allowbreak{} 10)} \\ \hline
\textit{Sala} & \texttt{CK\_\allowbreak{}Sala\_\allowbreak{}codigo} & \texttt{codigo LIKE \textquotesingle{}[A-HJ-NP-Z2-9]\allowbreak{}[A-HJ-NP-Z2-9]\allowbreak{}[A-HJ-NP-Z2-9]\allowbreak{}[A-HJ-NP-Z2-9]\allowbreak{}[A-HJ-NP-Z2-9]\allowbreak{}[A-HJ-NP-Z2-9]\allowbreak{}\textquotesingle{} COLLATE Latin1\_\allowbreak{}General\_\allowbreak{}100\_\allowbreak{}BIN2} \\ \hline
\textit{Sala} & \texttt{CK\_\allowbreak{}Sala\_\allowbreak{}ajustes} & \texttt{muros\_\allowbreak{}por\_\allowbreak{}jugador BETWEEN 1 AND 20 AND minutos\_\allowbreak{}reloj IN (3,\allowbreak{} 5,\allowbreak{} 10)} \\ \hline
\textit{Sala} & \texttt{CK\_\allowbreak{}Sala\_\allowbreak{}estado} & \texttt{estado IN (\textquotesingle{}ABIERTA\textquotesingle{},\allowbreak{}\textquotesingle{}EN\_\allowbreak{}PARTIDA\textquotesingle{},\allowbreak{}\textquotesingle{}CERRADA\textquotesingle{})} \\ \hline
\textit{Sala} & \texttt{CK\_\allowbreak{}Sala\_\allowbreak{}partida} & \texttt{estado <> \textquotesingle{}EN\_\allowbreak{}PARTIDA\textquotesingle{} OR id\_\allowbreak{}partida IS NOT NULL} \\ \hline
\textit{Sala\-Participante} & \texttt{CK\_\allowbreak{}Sala\allowbreak{}Participante\_\allowbreak{}plaza} & \texttt{plaza BETWEEN 1 AND 4} \\ \hline
\textit{Invitacion} & \texttt{CK\_\allowbreak{}Invitacion\_\allowbreak{}estado} & \texttt{estado IN (\textquotesingle{}PENDIENTE\textquotesingle{},\allowbreak{}\textquotesingle{}ACEPTADA\textquotesingle{},\allowbreak{}\textquotesingle{}RECHAZADA\textquotesingle{},\allowbreak{}\textquotesingle{}EXPIRADA\textquotesingle{})} \\ \hline
\textit{Invitacion} & \texttt{CK\_\allowbreak{}Invitacion\_\allowbreak{}distintos} & \texttt{id\_\allowbreak{}emisor <> id\_\allowbreak{}destinatario} \\ \hline
\textit{Mensaje} & \texttt{CK\_\allowbreak{}Mensaje\_\allowbreak{}canal} & \texttt{(canal IN (\textquotesingle{}PARTIDA\textquotesingle{},\allowbreak{}\textquotesingle{}ESPECTADORES\textquotesingle{}) AND id\_\allowbreak{}partida IS NOT NULL AND id\_\allowbreak{}sala IS NULL) OR (canal = \textquotesingle{}SALA\textquotesingle{} AND id\_\allowbreak{}sala IS NOT NULL AND id\_\allowbreak{}partida IS NULL) OR (canal IN (\textquotesingle{}GLOBAL\textquotesingle{},\allowbreak{}\textquotesingle{}AMIGOS\textquotesingle{}) AND id\_\allowbreak{}partida IS NULL AND id\_\allowbreak{}sala IS NULL)} \\ \hline
\textit{Mensaje} & \texttt{CK\_\allowbreak{}Mensaje\_\allowbreak{}texto} & \texttt{LEN(texto) > 0} \\ \hline
\multicolumn{3}{|l|}{\textbf{Clasificación y economía}} \\ \hline
\textit{Estadistica\-Modo} & \texttt{CK\_\allowbreak{}Estadistica\allowbreak{}Modo\_\allowbreak{}contadores} & \texttt{partidas\_\allowbreak{}ganadas + partidas\_\allowbreak{}perdidas <= partidas\_\allowbreak{}jugadas AND partidas\_\allowbreak{}ganadas >= 0 AND partidas\_\allowbreak{}perdidas >= 0 AND abandonos >= 0 AND tiempo\_\allowbreak{}total\_\allowbreak{}jugado >= 0 AND barreras\_\allowbreak{}colocadas >= 0 AND racha\_\allowbreak{}actual >= 0 AND mejor\_\allowbreak{}racha >= racha\_\allowbreak{}actual AND elo\_\allowbreak{}maximo >= puntos\_\allowbreak{}elo} \\ \hline
\textit{Historial\-Division} & \texttt{CK\_\allowbreak{}Historial\allowbreak{}Division\_\allowbreak{}cambio} & \texttt{id\_\allowbreak{}division\_\allowbreak{}anterior IS NULL OR id\_\allowbreak{}division\_\allowbreak{}anterior <> id\_\allowbreak{}division\_\allowbreak{}nueva} \\ \hline
\textit{Caja} & \texttt{CK\_\allowbreak{}Caja\_\allowbreak{}origen} & \texttt{origen IN (\textquotesingle{}COMPRA\textquotesingle{},\allowbreak{}\textquotesingle{}PARTIDA\textquotesingle{})} \\ \hline
\textit{Caja} & \texttt{CK\_\allowbreak{}Caja\_\allowbreak{}estado} & \texttt{estado IN (\textquotesingle{}SIN\_\allowbreak{}ABRIR\textquotesingle{},\allowbreak{}\textquotesingle{}ABIERTA\textquotesingle{},\allowbreak{}\textquotesingle{}CADUCADA\textquotesingle{})} \\ \hline
\textit{Caja} & \texttt{CK\_\allowbreak{}Caja\_\allowbreak{}caducidad} & \texttt{origen = \textquotesingle{}PARTIDA\textquotesingle{} OR fecha\_\allowbreak{}caducidad IS NULL} \\ \hline
\textit{Caja} & \texttt{CK\_\allowbreak{}Caja\_\allowbreak{}apertura} & \texttt{(estado = \textquotesingle{}ABIERTA\textquotesingle{} AND fecha\_\allowbreak{}apertura IS NOT NULL) OR (estado <> \textquotesingle{}ABIERTA\textquotesingle{} AND fecha\_\allowbreak{}apertura IS NULL)} \\ \hline
\textit{Caja\-Contenido} & \texttt{CK\_\allowbreak{}Caja\allowbreak{}Contenido\_\allowbreak{}rangos} & \texttt{numero BETWEEN 1 AND 3 AND (monedas\_\allowbreak{}conversion IS NULL OR monedas\_\allowbreak{}conversion > 0)} \\ \hline
\textit{Movimiento\-Moneda} & \texttt{CK\_\allowbreak{}Movimiento\allowbreak{}Moneda\_\allowbreak{}tipo} & \texttt{tipo IN (\textquotesingle{}PARTIDA\textquotesingle{},\allowbreak{}\textquotesingle{}NIVEL\textquotesingle{},\allowbreak{}\textquotesingle{}COMPRA\textquotesingle{},\allowbreak{}\textquotesingle{}COMPRA\_\allowbreak{}CAJA\textquotesingle{},\allowbreak{}\textquotesingle{}CONVERSION\textquotesingle{},\allowbreak{}\textquotesingle{}DEVOLUCION\textquotesingle{})} \\ \hline
\textit{Movimiento\-Moneda} & \texttt{CK\_\allowbreak{}Movimiento\allowbreak{}Moneda\_\allowbreak{}importe} & \texttt{(tipo IN (\textquotesingle{}COMPRA\textquotesingle{},\allowbreak{}\textquotesingle{}COMPRA\_\allowbreak{}CAJA\textquotesingle{}) AND importe < 0) OR (tipo IN (\textquotesingle{}PARTIDA\textquotesingle{},\allowbreak{}\textquotesingle{}NIVEL\textquotesingle{},\allowbreak{}\textquotesingle{}CONVERSION\textquotesingle{}) AND importe > 0) OR (tipo = \textquotesingle{}DEVOLUCION\textquotesingle{} AND importe <> 0)} \\ \hline
\textit{Movimiento\-Moneda} & \texttt{CK\_\allowbreak{}Movimiento\allowbreak{}Moneda\_\allowbreak{}origen} & \texttt{(tipo = \textquotesingle{}PARTIDA\textquotesingle{} AND id\_\allowbreak{}partida IS NOT NULL) OR (tipo = \textquotesingle{}COMPRA\textquotesingle{} AND id\_\allowbreak{}objeto IS NOT NULL) OR (tipo = \textquotesingle{}COMPRA\_\allowbreak{}CAJA\textquotesingle{} AND id\_\allowbreak{}caja IS NOT NULL) OR (tipo = \textquotesingle{}CONVERSION\textquotesingle{} AND id\_\allowbreak{}caja IS NOT NULL AND id\_\allowbreak{}objeto IS NOT NULL) OR tipo IN (\textquotesingle{}NIVEL\textquotesingle{},\allowbreak{}\textquotesingle{}DEVOLUCION\textquotesingle{})} \\ \hline
\multicolumn{3}{|l|}{\textbf{Relaciones entre jugadores y tutorial}} \\ \hline
\textit{Amistad} & \texttt{CK\_\allowbreak{}Amistad\_\allowbreak{}par\_\allowbreak{}ordenado} & \texttt{id\_\allowbreak{}usuario\_\allowbreak{}a < id\_\allowbreak{}usuario\_\allowbreak{}b} \\ \hline
\textit{Solicitud} & \texttt{CK\_\allowbreak{}Solicitud\_\allowbreak{}distintos} & \texttt{id\_\allowbreak{}solicitante <> id\_\allowbreak{}destinatario} \\ \hline
\textit{Silencio} & \texttt{CK\_\allowbreak{}Silencio\_\allowbreak{}distintos} & \texttt{id\_\allowbreak{}usuario <> id\_\allowbreak{}silenciado} \\ \hline
\textit{Bloqueo} & \texttt{CK\_\allowbreak{}Bloqueo\_\allowbreak{}distintos} & \texttt{id\_\allowbreak{}bloqueador <> id\_\allowbreak{}bloqueado} \\ \hline
\textit{Progreso\-Tutorial} & \texttt{CK\_\allowbreak{}Progreso\allowbreak{}Tutorial\_\allowbreak{}paso} & \texttt{paso\_\allowbreak{}actual >= 1} \\ \hline
\multicolumn{3}{|l|}{\textbf{Moderación y bitácoras}} \\ \hline
\textit{Reporte} & \texttt{CK\_\allowbreak{}Reporte\_\allowbreak{}distintos} & \texttt{id\_\allowbreak{}denunciante <> id\_\allowbreak{}reportado} \\ \hline
\textit{Reporte} & \texttt{CK\_\allowbreak{}Reporte\_\allowbreak{}estado} & \texttt{(estado = \textquotesingle{}PENDIENTE\textquotesingle{} AND id\_\allowbreak{}moderador IS NULL AND fecha\_\allowbreak{}asignacion IS NULL AND fecha\_\allowbreak{}resolucion IS NULL) OR (estado = \textquotesingle{}EN\_\allowbreak{}REVISION\textquotesingle{} AND id\_\allowbreak{}moderador IS NOT NULL AND fecha\_\allowbreak{}asignacion IS NOT NULL AND fecha\_\allowbreak{}resolucion IS NULL) OR (estado IN (\textquotesingle{}RESUELTO\_\allowbreak{}SIN\_\allowbreak{}SANCION\textquotesingle{},\allowbreak{}\textquotesingle{}RESUELTO\_\allowbreak{}CON\_\allowbreak{}SANCION\textquotesingle{}) AND id\_\allowbreak{}moderador IS NOT NULL AND nota\_\allowbreak{}resolucion IS NOT NULL AND fecha\_\allowbreak{}resolucion IS NOT NULL)} \\ \hline
\textit{Sancion} & \texttt{CK\_\allowbreak{}Sancion\_\allowbreak{}ambito} & \texttt{ambito IN (\textquotesingle{}CUENTA\textquotesingle{},\allowbreak{}\textquotesingle{}CHAT\textquotesingle{})} \\ \hline
\textit{Sancion} & \texttt{CK\_\allowbreak{}Sancion\_\allowbreak{}tipo} & \texttt{(tipo = \textquotesingle{}TEMPORAL\textquotesingle{} AND fecha\_\allowbreak{}fin IS NOT NULL AND fecha\_\allowbreak{}fin > fecha\_\allowbreak{}inicio) OR (tipo = \textquotesingle{}PERMANENTE\textquotesingle{} AND fecha\_\allowbreak{}fin IS NULL)} \\ \hline
\textit{Sancion} & \texttt{CK\_\allowbreak{}Sancion\_\allowbreak{}sustitucion} & \texttt{id\_\allowbreak{}sancion\_\allowbreak{}sustituta IS NULL OR fecha\_\allowbreak{}retiro IS NOT NULL} \\ \hline
\textit{Sancion} & \texttt{CK\_\allowbreak{}Sancion\_\allowbreak{}distintos} & \texttt{id\_\allowbreak{}usuario <> id\_\allowbreak{}moderador} \\ \hline
\textit{Apelacion} & \texttt{CK\_\allowbreak{}Apelacion\_\allowbreak{}estado} & \texttt{(estado = \textquotesingle{}PENDIENTE\textquotesingle{} AND id\_\allowbreak{}moderador IS NULL AND fecha\_\allowbreak{}asignacion IS NULL AND fecha\_\allowbreak{}resolucion IS NULL) OR (estado = \textquotesingle{}EN\_\allowbreak{}REVISION\textquotesingle{} AND id\_\allowbreak{}moderador IS NOT NULL AND fecha\_\allowbreak{}asignacion IS NOT NULL AND fecha\_\allowbreak{}resolucion IS NULL) OR (estado IN (\textquotesingle{}ACEPTADA\textquotesingle{},\allowbreak{}\textquotesingle{}RECHAZADA\textquotesingle{}) AND id\_\allowbreak{}moderador IS NOT NULL AND nota\_\allowbreak{}resolucion IS NOT NULL AND fecha\_\allowbreak{}resolucion IS NOT NULL)} \\ \hline
\textit{Bitacora\-Moderacion} & \texttt{CK\_\allowbreak{}Bitacora\allowbreak{}Moderacion\_\allowbreak{}accion} & \texttt{accion IN (\textquotesingle{}REPORTE\_\allowbreak{}TOMADO\textquotesingle{},\allowbreak{}\textquotesingle{}REPORTE\_\allowbreak{}LIBERADO\textquotesingle{},\allowbreak{}\textquotesingle{}REPORTE\_\allowbreak{}RESUELTO\textquotesingle{},\allowbreak{}\textquotesingle{}SANCION\_\allowbreak{}APLICADA\textquotesingle{} \textquotesingle{}APELACION\_\allowbreak{}RESUELTA\textquotesingle{},\allowbreak{}\textquotesingle{}ROL\_\allowbreak{}OTORGADO\textquotesingle{},\allowbreak{}\textquotesingle{}ROL\_\allowbreak{}RETIRADO\textquotesingle{},\allowbreak{}\textquotesingle{}CONSULTA\_\allowbreak{}BITACORA\textquotesingle{},\allowbreak{}\textquotesingle{}INTENTO\_\allowbreak{}SIN\_\allowbreak{}PERMISO\textquotesingle{})} \\ \hline
\textit{Bitacora\-Acceso} & \texttt{CK\_\allowbreak{}Bitacora\allowbreak{}Acceso\_\allowbreak{}resultado} & \texttt{resultado IN (\textquotesingle{}EXITO\textquotesingle{},\allowbreak{}\textquotesingle{}USUARIO\_\allowbreak{}INEXISTENTE\textquotesingle{},\allowbreak{}\textquotesingle{}BLOQUEO\_\allowbreak{}VIGENTE\textquotesingle{},\allowbreak{}\textquotesingle{}CONTRASENA\_\allowbreak{}INCORRECTA\textquotesingle{},\allowbreak{}\textquotesingle{}CUENTA\_\allowbreak{}PENDIENTE\textquotesingle{},\allowbreak{}\textquotesingle{}CUENTA\_\allowbreak{}SANCIONADA\textquotesingle{} \textquotesingle{}SEGUNDO\_\allowbreak{}FACTOR\_\allowbreak{}FALLIDO\textquotesingle{},\allowbreak{}\textquotesingle{}ABANDONADO\textquotesingle{},\allowbreak{}\textquotesingle{}SUSPENSION\_\allowbreak{}VENCIDA\textquotesingle{},\allowbreak{}\textquotesingle{}REGISTRO\_\allowbreak{}EXITOSO\textquotesingle{},\allowbreak{}\textquotesingle{}REGISTRO\_\allowbreak{}RECHAZADO\textquotesingle{},\allowbreak{}\textquotesingle{}REGISTRO\_\allowbreak{}SIN\_\allowbreak{}CORREO\textquotesingle{},\allowbreak{} \textquotesingle{}ALTA\_\allowbreak{}INVITADO\textquotesingle{},\allowbreak{}\textquotesingle{}VINCULACION\_\allowbreak{}EXITOSA\textquotesingle{},\allowbreak{}\textquotesingle{}CUENTA\_\allowbreak{}VERIFICADA\textquotesingle{},\allowbreak{}\textquotesingle{}RECUPERACION\_\allowbreak{}SOLICITADA\textquotesingle{},\allowbreak{}\textquotesingle{}RECUPERACION\_\allowbreak{}CORREO\_\allowbreak{}INEXISTENTE\textquotesingle{},\allowbreak{} \textquotesingle{}RECUPERACION\_\allowbreak{}EXITOSA\textquotesingle{},\allowbreak{}\textquotesingle{}CIERRE\_\allowbreak{}VOLUNTARIO\textquotesingle{},\allowbreak{}\textquotesingle{}CIERRE\_\allowbreak{}REMOTO\textquotesingle{},\allowbreak{}\textquotesingle{}CAMBIO\_\allowbreak{}CONTRASENA\textquotesingle{},\allowbreak{}\textquotesingle{}CAMBIO\_\allowbreak{}CORREO\_\allowbreak{}SOLICITADO\textquotesingle{},\allowbreak{} \textquotesingle{}CAMBIO\_\allowbreak{}CORREO\_\allowbreak{}CONFIRMADO\textquotesingle{},\allowbreak{}\textquotesingle{}SEGUNDO\_\allowbreak{}FACTOR\_\allowbreak{}ACTIVADO\textquotesingle{},\allowbreak{}\textquotesingle{}SEGUNDO\_\allowbreak{}FACTOR\_\allowbreak{}DESACTIVADO\textquotesingle{},\allowbreak{}\textquotesingle{}CUENTA\_\allowbreak{}ELIMINADA\textquotesingle{})} \\ \hline
\end{longtable}
\end{center}


\subsubsection*{\texttt{DEFAULT}}
\addcontentsline{toc}{subsubsection}{DEFAULT}

\bdNumOmision{} columnas tienen un valor por omisión: el estado inicial de la fila, como \at{PENDIENTE} o \at{SIN\_ABRIR}; los contadores en cero y el elo en 1000, con los que nace una cuenta (\mbox{CU-02} \mbox{RN-06}); los indicadores en su valor habitual; y la fecha actual, \at{SYSUTCDATETIME()}, en las fechas que registran cuándo se creó la fila. Así, un \at{INSERT} solo tiene que dar lo que depende del caso.

% Archivo generado por bd/generar_tablas.ps1 a partir de bd/crear_tablas.sql. No editar a mano.
\begin{center}\small
\begin{longtable}{|L{0.22\textwidth}|L{0.68\textwidth}|}
\hline
\textbf{Entidad} & \textbf{Columna y valor por omisión} \\ \hline
\endfirsthead
\hline
\textbf{Entidad} & \textbf{Columna y valor por omisión} \\ \hline
\endhead
\textit{Objeto\-Cosmetico} & \texttt{precio} = \texttt{0}; \texttt{nivel\_\allowbreak{}requerido} = \texttt{1}; \texttt{a\_\allowbreak{}la\_\allowbreak{}venta} = \texttt{0}; \texttt{activo} = \texttt{1}; \texttt{es\_\allowbreak{}inicial} = \texttt{0}; \texttt{es\_\allowbreak{}predeterminado} = \texttt{0} \\ \hline
\textit{Modo} & \texttt{activo} = \texttt{1} \\ \hline
\textit{Tipo\-Caja} & \texttt{activo} = \texttt{1} \\ \hline
\textit{Usuario} & \texttt{estado\_\allowbreak{}cuenta} = \texttt{\textquotesingle{}PENDIENTE\textquotesingle{}}; \texttt{rol} = \texttt{\textquotesingle{}JUGADOR\textquotesingle{}}; \texttt{idioma\_\allowbreak{}preferido} = \texttt{\textquotesingle{}es-MX\textquotesingle{}}; \texttt{id\_\allowbreak{}icono} = \texttt{1}; \texttt{permite\_\allowbreak{}espectadores} = \texttt{1}; \texttt{nivel} = \texttt{1}; \texttt{experiencia} = \texttt{0}; \texttt{saldo\_\allowbreak{}monedas} = \texttt{0}; \texttt{cajas\_\allowbreak{}sin\_\allowbreak{}raro} = \texttt{0}; \texttt{intentos\_\allowbreak{}fallidos} = \texttt{0}; \texttt{doble\_\allowbreak{}factor\_\allowbreak{}habilitado} = \texttt{0}; \texttt{fecha\_\allowbreak{}registro} = \texttt{SYSUTCDATETIME()} \\ \hline
\textit{Sesion} & \texttt{fecha\_\allowbreak{}inicio} = \texttt{SYSUTCDATETIME()}; \texttt{fecha\_\allowbreak{}ultimo\_\allowbreak{}uso} = \texttt{SYSUTCDATETIME()} \\ \hline
\textit{Codigo\-Segundo\-Factor} & \texttt{fecha\_\allowbreak{}generacion} = \texttt{SYSUTCDATETIME()}; \texttt{intentos} = \texttt{0}; \texttt{estado} = \texttt{\textquotesingle{}PENDIENTE\textquotesingle{}} \\ \hline
\textit{Token\-Verificacion\-Correo} & \texttt{fecha\_\allowbreak{}generacion} = \texttt{SYSUTCDATETIME()}; \texttt{estado} = \texttt{\textquotesingle{}PENDIENTE\textquotesingle{}} \\ \hline
\textit{Token\-Recuperacion} & \texttt{fecha\_\allowbreak{}generacion} = \texttt{SYSUTCDATETIME()}; \texttt{intentos} = \texttt{0}; \texttt{estado} = \texttt{\textquotesingle{}PENDIENTE\textquotesingle{}} \\ \hline
\textit{Aceptacion\-Terminos} & \texttt{fecha\_\allowbreak{}aceptacion} = \texttt{SYSUTCDATETIME()} \\ \hline
\textit{Historial\-Nickname} & \texttt{fecha\_\allowbreak{}cambio} = \texttt{SYSUTCDATETIME()} \\ \hline
\textit{Avatar} & \texttt{fecha\_\allowbreak{}subida} = \texttt{SYSUTCDATETIME()}; \texttt{vigente} = \texttt{1} \\ \hline
\textit{Usuario\-Objeto} & \texttt{fecha\_\allowbreak{}obtencion} = \texttt{SYSUTCDATETIME()} \\ \hline
\textit{Partida} & \texttt{estado} = \texttt{\textquotesingle{}EN\_\allowbreak{}CURSO\textquotesingle{}}; \texttt{fecha\_\allowbreak{}inicio} = \texttt{SYSUTCDATETIME()}; \texttt{turno\_\allowbreak{}actual} = \texttt{1} \\ \hline
\textit{Participacion} & \texttt{conectado} = \texttt{1} \\ \hline
\textit{Jugada} & \texttt{fecha\_\allowbreak{}jugada} = \texttt{SYSUTCDATETIME()}; \texttt{deshecha} = \texttt{0} \\ \hline
\textit{Oferta\-Tablas} & \texttt{fecha\_\allowbreak{}oferta} = \texttt{SYSUTCDATETIME()}; \texttt{estado} = \texttt{\textquotesingle{}PENDIENTE\textquotesingle{}} \\ \hline
\textit{Enlace\-Espectador} & \texttt{fecha\_\allowbreak{}creacion} = \texttt{SYSUTCDATETIME()} \\ \hline
\textit{Cola\-Emparejamiento} & \texttt{fecha\_\allowbreak{}entrada} = \texttt{SYSUTCDATETIME()} \\ \hline
\textit{Sala} & \texttt{permite\_\allowbreak{}espectadores} = \texttt{1}; \texttt{estado} = \texttt{\textquotesingle{}ABIERTA\textquotesingle{}}; \texttt{fecha\_\allowbreak{}creacion} = \texttt{SYSUTCDATETIME()}; \texttt{fecha\_\allowbreak{}ultima\_\allowbreak{}actividad} = \texttt{SYSUTCDATETIME()} \\ \hline
\textit{Sala\-Participante} & \texttt{listo} = \texttt{0}; \texttt{fecha\_\allowbreak{}union} = \texttt{SYSUTCDATETIME()} \\ \hline
\textit{Invitacion} & \texttt{fecha\_\allowbreak{}invitacion} = \texttt{SYSUTCDATETIME()}; \texttt{estado} = \texttt{\textquotesingle{}PENDIENTE\textquotesingle{}} \\ \hline
\textit{Mensaje} & \texttt{fecha\_\allowbreak{}envio} = \texttt{SYSUTCDATETIME()} \\ \hline
\textit{Estadistica\-Modo} & \texttt{puntos\_\allowbreak{}elo} = \texttt{1000}; \texttt{elo\_\allowbreak{}maximo} = \texttt{1000}; \texttt{partidas\_\allowbreak{}jugadas} = \texttt{0}; \texttt{partidas\_\allowbreak{}ganadas} = \texttt{0}; \texttt{partidas\_\allowbreak{}perdidas} = \texttt{0}; \texttt{racha\_\allowbreak{}actual} = \texttt{0}; \texttt{mejor\_\allowbreak{}racha} = \texttt{0}; \texttt{tiempo\_\allowbreak{}total\_\allowbreak{}jugado} = \texttt{0}; \texttt{barreras\_\allowbreak{}colocadas} = \texttt{0}; \texttt{abandonos} = \texttt{0} \\ \hline
\textit{Historial\-Division} & \texttt{fecha\_\allowbreak{}cambio} = \texttt{SYSUTCDATETIME()} \\ \hline
\textit{Caja} & \texttt{estado} = \texttt{\textquotesingle{}SIN\_\allowbreak{}ABRIR\textquotesingle{}}; \texttt{fecha\_\allowbreak{}obtencion} = \texttt{SYSUTCDATETIME()} \\ \hline
\textit{Movimiento\-Moneda} & \texttt{fecha\_\allowbreak{}movimiento} = \texttt{SYSUTCDATETIME()} \\ \hline
\textit{Amistad} & \texttt{fecha\_\allowbreak{}amistad} = \texttt{SYSUTCDATETIME()} \\ \hline
\textit{Solicitud} & \texttt{fecha\_\allowbreak{}solicitud} = \texttt{SYSUTCDATETIME()} \\ \hline
\textit{Silencio} & \texttt{fecha\_\allowbreak{}silencio} = \texttt{SYSUTCDATETIME()} \\ \hline
\textit{Bloqueo} & \texttt{fecha\_\allowbreak{}bloqueo} = \texttt{SYSUTCDATETIME()} \\ \hline
\textit{Progreso\-Tutorial} & \texttt{fecha\_\allowbreak{}actualizacion} = \texttt{SYSUTCDATETIME()} \\ \hline
\textit{Reporte} & \texttt{fecha\_\allowbreak{}reporte} = \texttt{SYSUTCDATETIME()}; \texttt{estado} = \texttt{\textquotesingle{}PENDIENTE\textquotesingle{}} \\ \hline
\textit{Sancion} & \texttt{fecha\_\allowbreak{}inicio} = \texttt{SYSUTCDATETIME()} \\ \hline
\textit{Apelacion} & \texttt{marca\_\allowbreak{}lenguaje} = \texttt{0}; \texttt{fecha\_\allowbreak{}apelacion} = \texttt{SYSUTCDATETIME()}; \texttt{estado} = \texttt{\textquotesingle{}PENDIENTE\textquotesingle{}} \\ \hline
\textit{Bitacora\-Moderacion} & \texttt{fecha\_\allowbreak{}hora} = \texttt{SYSUTCDATETIME()} \\ \hline
\textit{Bitacora\-Acceso} & \texttt{fecha\_\allowbreak{}hora} = \texttt{SYSUTCDATETIME()} \\ \hline
\end{longtable}
\end{center}


\subsubsection*{Columnas calculadas e índices}
\addcontentsline{toc}{subsubsection}{Columnas calculadas e índices}

Las \bdNumCalculadas{} columnas calculadas, \at{diferencia\_elo}, \at{terminada}, \at{convertido}, \at{completada} y \at{activa}, se declaran con \at{AS (expresión)} y no se almacenan: se obtienen de otras columnas de la misma fila cada vez que se leen, así que no pueden contradecirlas. Por último, \bdNumIndices{} índices no únicos aceleran las consultas frecuentes que no siguen la llave primaria: el ranking por modo, la búsqueda en la cola de emparejamiento, el chat de cada canal, las sesiones abiertas de una cuenta, la cola de reportes, las sanciones vigentes y la bitácora por fecha.
